\documentclass[a4paper,12pt]{article}

% 基础宏包配置
\usepackage{ctex}      % 中文支持
\usepackage{geometry}  % 页面布局
\usepackage{booktabs}  % 专业表格线
\usepackage{longtable} % 长表格
\usepackage{array}     % 表格列控制
\usepackage{enumitem}  % 列表定制
\usepackage{titlesec}  % 标题格式
\usepackage{fancyhdr}  % 页眉页脚
\usepackage{xcolor}    % 颜色支持

% 页面边距设置
\geometry{left=2.5cm, right=2.5cm, top=3cm, bottom=3cm}

% 颜色定义 - 采用冷峻的深色系以匹配学术风格
\definecolor{academicBlue}{RGB}{0, 40, 80}
\definecolor{formalGray}{RGB}{60, 60, 60}

% 标题格式定制
\titleformat{\section}
{\normalfont\Large\bfseries\color{academicBlue}}{\thesection}{1em}{}
\titleformat{\subsection}
{\normalfont\large\bfseries\color{formalGray}}{\thesubsection}{1em}{}
\titleformat{\subsubsection}
{\normalfont\normalsize\bfseries}{\thesubsubsection}{1em}{}

% 列表样式定制
\setlist[itemize]{label=$\bullet$, leftmargin=*, itemsep=2pt, topsep=2pt}

% 页眉页脚
\pagestyle{fancy}
\fancyhf{}
\lhead{\small \textbf{2024年度环境、社会及治理（ESG）报告}}
\rhead{\small 管理层讨论与分析}
\cfoot{\thepage}

\title{\textbf{2024年度ESG报告}\\ \large 可持续发展管理体系与绩效综述}
\author{本公司董事会}
\date{2025年3月}

\begin{document}

\maketitle

\section{治理维度：公司治理架构与经济价值创造}

在2024财政年度内，鉴于全球宏观经济环境的不确定性加剧以及汽车产业向“新质生产力”转型的结构性压力，本公司坚持执行“三直四化”（即电动化、智能网联化、高端化及国际化）的战略部署。通过对公司治理机制的持续优化以及对经营风险的系统性管控，实现了经济效益的稳健增长与股东回报机制的有效落实。

\subsection{经济绩效与资本回报}

报告期内，得益于海外市场的渗透率提升以及产品结构的高端化调整，公司的主要财务指标呈现出显著的增长态势。关于利润分配政策，基于对现金流状况的审慎评估及对长期价值投资理念的遵循，年度内实施了双重现金分红方案。

下表列示了经审计的2024年度核心经济绩效指标：

\begin{table}[h]
\centering
\caption{2024年度核心经济绩效指标表}
\renewcommand{\arraystretch}{1.3}
\begin{tabular}{p{0.4\textwidth} p{0.3\textwidth} p{0.2\textwidth}}
\toprule
\textbf{指标名称} & \textbf{数值} & \textbf{变动情况} \\
\midrule
营业收入 & \textbf{372.18 亿元} & +37.63\% \\
归属于母公司股东的净利润 & \textbf{41.16 亿元} & 增长 \\
现金股利分派总额 & \textbf{44.28 亿元} & - \\
分红实施频次 & \textbf{2 次} & - \\
\bottomrule
\end{tabular}
\end{table}

\subsection{合规体系与供应链尽职管理}

为确保企业运营的合规性与透明度，公司建立并维持了一套符合国际标准的合规管理体系。同时，针对供应链的ESG风险，实施了严格的准入与审核机制。

\begin{itemize}
    \item \textbf{信息披露合规性}：
    鉴于内部控制体系的有效运行，公司已连续第13个年度被证券交易所授予信息披露工作“A级”评价，该评价反映了监管机构对公司透明度的认可。
    \item \textbf{合规体系认证}：
    报告期内，公司完成了ISO 37301合规管理体系的认证程序，标志着内部合规管理架构已满足国际标准化的要求。
    \item \textbf{供应链管理}：
    \begin{itemize}
        \item 截至报告期末，合格供应商名录共包含 \textbf{530 家} 企业，其中河南本地供应商占比约为 \textbf{20\%}，体现了对区域产业集群发展的支持。
        \item 为规避商业道德风险，已与超过 \textbf{900 家} 供应商（含潜在合作方）签署了具有法律约束力的廉洁协议。
    \end{itemize}
\end{itemize}

\section{社会维度：技术资本积累与相关方权益}

社会责任的履行被视为公司长期战略的重要组成部分，其核心在于通过高强度的研发投入推动行业技术进步，并通过完善的人力资源管理体系与公益投入机制，实现社会价值的最大化。

\subsection{研发创新与知识产权管理}

作为技术驱动型企业，研发资源的配置被优先用于新能源与智能网联核心技术的突破。报告期内，研发活动的投入强度与产出效率均维持在行业领先水平。

\begin{table}[h]
\centering
\caption{2024年度研发资源投入与产出统计}
\renewcommand{\arraystretch}{1.3}
\begin{tabular}{l l}
\toprule
\textbf{项目类别} & \textbf{统计数据} \\
\midrule
研发费用投入 & 17.88 亿元（占营业收入比重 4.80\%） \\
研发人员配置 & 3,505 人（其中拥有博士学位人员 16 人） \\
新增授权专利 & 176 项（其中发明专利 114 项） \\
标准制定参与 & 317 项（累计数） \\
\bottomrule
\end{tabular}
\end{table}

\subsubsection{核心技术体系}
本年度内的技术创新主要集中于以下关键领域，旨在通过技术迭代解决行业痛点：

\begin{itemize}
    \item \textbf{YOS车机操作系统}：
    开发了基于软硬解耦架构的专用操作系统，该系统具备全栈式空中下载（OTA）升级能力，确立了软件定义汽车的技术基础。
    \item \textbf{电子电气架构（C架构）}：
    实施了以“中央计算+区域控制”为特征的新型架构，该架构通过高度集成化设计，实现了线束复杂度的显著降低与整车轻量化。
    \item \textbf{电驱动系统效率优化}：
    在多合一动力域控制器中导入碳化硅（SiC）功率模块技术。相较于传统硅基IGBT模块，该技术方案使开关损耗降低了 \textbf{50\%-70\%}，从而提升了系统的综合能效。
    \item \textbf{热管理技术}：
    研发并应用了多源超低温热泵系统，该系统通过对环境热能及电机废热的回收利用，在低温工况下实现了约 \textbf{30\%} 的采暖能耗降低。
    \item \textbf{智能网联运营}：
    L4级自动驾驶巴士的累计安全运营里程已超过 \textbf{1,700 万公里}，且公司已被纳入国家首批智能网联汽车准入试点名单。
\end{itemize}

\subsection{人力资本与职业健康安全}
人力资源被视为公司核心资产，公司致力于构建多元包容的职场环境，并依据ISO 45001标准执行职业健康安全管理。

\begin{itemize}
    \item \textbf{人员结构}：截至2024年12月31日，员工总数为 16,707 人，其中女性员工占比为 10.02\%。
    \item \textbf{安全投入与绩效}：
    \begin{itemize}
        \item 劳动防护用品专项投入：\textbf{5,500 余万元}。
        \item 生产环境改善专项投入：\textbf{2,200 余万元}。
        \item 安全教育培训专项投入：\textbf{180 余万元}。
        \item 职业健康绩效：报告期内 \textbf{无新增职业病} 病例记录。
    \end{itemize}
\end{itemize}

\subsection{企业公民责任履行}
公司通过战略性公益项目的实施，积极回应社会关切。2024年度，对外捐赠及公益项目支出总额超过 \textbf{4,000 万元}。

\begin{table}[h]
\centering
\caption{2024年度主要社会公益项目实施情况}
\begin{tabular}{p{0.3\textwidth} p{0.6\textwidth}}
\toprule
\textbf{项目标识} & \textbf{实施内容与社会效益} \\
\midrule
儿童交通安全公益行 & 该项目覆盖全国13个省市，通过模拟盲区体验等教育手段，旨在降低近万名受众儿童的交通安全风险。 \\
全球零碳森林计划 & 在南美某国实施生态修复项目，累计种植耐旱树种36,000余棵，旨在通过碳汇机制抵消运营碳足迹。 \\
\bottomrule
\end{tabular}
\end{table}

\section{环境维度：环境管理体系与绩效指标}

公司承诺遵守所有适用的环境法律法规，并将绿色低碳原则纳入产品全生命周期管理。以下数据基于公司内部环境监测系统及第三方核查结果编制。

\subsection{温室气体排放与能源消耗}
为量化评估气候变化风险及应对成效，公司对关键环境指标进行了系统性核算。

\begin{table}[h]
\centering
\caption{2024年度环境绩效指标摘要}
\renewcommand{\arraystretch}{1.3}
\begin{tabular}{p{0.5\textwidth} p{0.4\textwidth}}
\toprule
\textbf{指标项目} & \textbf{核算数值} \\
\midrule
范围1排放量（直接温室气体排放） & 57,330.53 tCO$_2$e \\
范围2排放量（能源间接温室气体排放） & 122,786.63 tCO$_2$e \\
\textbf{温室气体排放强度} & \textbf{0.07 tCO$_2$e/万元营收} \\
综合能源消费量 & 34,560.94 吨标准煤 \\
\bottomrule
\end{tabular}
\end{table}

\subsection{节能减排措施实施情况}
针对上述排放源，公司实施了以下技术改造与管理优化措施：

\begin{itemize}
    \item \textbf{能效提升项目}：
    报告期内，完成了7项主要节能技术改造工程。通过设备更新与工艺优化，实现了年均 \textbf{1,258 MWh} 的电力节约及 \textbf{11.34 万立方米} 的天然气节约。
    \item \textbf{余热回收利用}：
    通过工业余热回收系统的部署，报告期内回收利用热能总量约为 \textbf{29,306 GJ}，有效降低了化石能源的直接消耗。
    \item \textbf{清洁能源替代}：
    \begin{itemize}
        \item 分布式光伏发电量：\textbf{155.68 万 kWh}（折合减碳约 835 吨）。
        \item 绿色电力采购量：\textbf{13,580 MWh}。
    \end{itemize}
\end{itemize}

\subsection{资源循环与废弃物管理}
资源利用效率的提升与废弃物的合规处置是环境管理体系的重点内容。

\begin{table}[h]
\centering
\caption{2024年度资源利用与废弃物处置统计}
\begin{tabular}{l l}
\toprule
\textbf{类别} & \textbf{统计指标} \\
\midrule
\textbf{水资源管理} & \\
- 新水取用量 & 141.25 万吨 \\
- 循环用水量 & 6,465.36 万吨 \\
- 工业用水重复利用率 & \textbf{97.8\%} \\
\midrule
\textbf{固体废弃物} & \\
- 一般工业固废资源化利用量 & 61,485.76 吨 \\
- 危险废物产生量 & 4,110.84 吨（合规处置率 100\%） \\
\midrule
\textbf{绿色包装} & \\
- 国产零部件周转包装占比 & \textbf{89\%} \\
\bottomrule
\end{tabular}
\end{table}

\end{document}